\chapter*{Conclusion générale}
\phantomsection
\addcontentsline{toc}{chapter}{Conclusion générale}

Ce projet a étudié la gravité des blessures parmi les usagers des accidents corporels
enregistrés en France métropolitaine entre 2019 et 2024. Il a suivi la chaîne complète
d'une étude de données~: compréhension des sources, audit, préparation, analyse
statistique, modélisation et évaluation sur une année ultérieure.

Le premier résultat est une table d'analyse traçable~: 702\,823 observations d'usagers
pour 309\,756 accidents distincts. Sa construction a demandé de résoudre les différences
de schéma, les codes spéciaux, les identifiants, les valeurs atypiques et les lieux
multiples. Les règles retenues préservent les usagers, signalent les désaccords et
maintiennent les valeurs sources nécessaires à la vérification. Les 365 gravités
inconnues restent identifiées, puis sont exclues des calculs supervisés.

Les analyses de 2019-2022 répondent à la question des différences observées. Les
proportions de décès ou d'hospitalisations sont plus élevées chez les piétons que chez
les conducteurs, chez les personnes de 75 ans ou plus que chez les 35-44 ans, de nuit
sans éclairage public qu'en plein jour, et hors agglomération qu'en agglomération.
Le bootstrap par accident rend compte d'une variabilité en respectant le regroupement
des usagers. La lecture annuelle et la sensibilité aux décès seuls apportent des
compléments, sans effacer les limites de collecte ni transformer ces associations
en effets causaux.

La partie prédictive apporte une réponse complémentaire. Avec onze variables et une
régression logistique régularisée, les probabilités sont mieux estimées que par la
seule fréquence de l'apprentissage. Le réglage choisi sur 2023 se maintient sur
2024~: la log-loss diminue d'environ 14,93\,\% par rapport à la référence, avec une
ROC-AUC proche de 0,770. Le projet établit ainsi un apport mesurable dans le cadre
de l'évaluation retenue.

Les diagnostics précisent la portée de ce résultat. Au seuil de 0,5, le rappel est
limité à 18,22\,\%. Les probabilités élevées sont souvent surestimées et le classement
est moins performant chez les piétons. L'absence d'âge ou de sexe porte également un
signal de saisie particulier.

Le chapitre~\ref{chap:comparaison} reprend ces trois réserves. Un modèle d'arbres
boostés, évalué avec les mêmes variables et le même protocole, améliore les quatre
scores~: la log-loss descend à 0,3845 et la ROC-AUC atteint 0,789 sur 2024. Le gain
porte surtout sur le haut du classement, là où le modèle additif plafonne faute de
pouvoir représenter des combinaisons de caractéristiques. La régression logistique
conserve en revanche son avantage de lisibilité. L'arbitrage précision-rappel, rendu
explicite, montre qu'un rappel de 50\,\% est atteignable en signalant un usager sur
cinq, avec une règle fixée sur 2023 qui se transporte à 2024. Enfin, le signal porté
par l'absence d'âge est daté de 2021, rapproché d'un changement de collecte documenté
et chiffré~: il représente environ 0,6 point de ROC-AUC.

Cette dernière comparaison est intervenue après consultation des scores de 2024.
Ses réglages ont été choisis sur 2023 seule, mais une année réellement vierge serait
nécessaire pour trancher définitivement entre les deux familles de modèles.

La principale contribution du projet tient à la cohérence entre la question,
l'unité d'analyse, les transformations et l'évaluation. Les résultats restent
conditionnels à la présence dans BAAC. Ils ne décrivent pas un risque par trajet,
ne permettent pas une interprétation causale des coefficients et ne démontrent
pas l'intérêt d'un outil de décision en temps réel. Cette délimitation permet
de défendre précisément ce qui a été réalisé.

Plusieurs prolongements restent possibles~: approfondir la calibration, étudier plus
finement les mécanismes d'information manquante, introduire les champs de lieux écartés
du premier modèle, ou évaluer les deux familles sur un millésime postérieur à 2024.
Ce dernier point est le plus utile, puisqu'il fournirait le test indépendant qui manque
à la comparaison des modèles. Pour mesurer un risque par déplacement ou éclairer une
décision d'aménagement, des données d'exposition et un objectif d'usage seraient
également nécessaires.
